% !TeX root = ../se200_identity_regimes.tex

\section{Formal Setting}
\label{sec:formal_setting}

This section introduces the structural objects over which regime profiles are defined.
We first identify the types of structures that may carry identity,
and then specify the transformation families under which identity must remain stable.

%--------------------------------------------
\subsection{Identity Carriers}
\label{subsec:identity_carriers}

This subsection introduces the structural objects
that may bear identity under regime profiles.

\begin{definition}[Identity carrier]
\label{def:identity_carrier}
An identity carrier is a structural object in $\mathcal{R}$
whose identity is determined under a regime profile.
\end{definition}


Identity carriers are structural positions within $\mathcal{R}$,
not ontological types.
Regime assignment depends solely on transformation behavior,
not carrier category.
Identity may be borne by different structural forms,
each supporting distinct identity and persistence behavior under transformation.


\begin{example}[Carrier types]
\label{ex:carrier_types}
Identity carriers include structural forms such as:
\begin{itemize}
\item enduring non-normative referents,
\item normative structures,
\item time-indexed occurrences,
\item applicability contexts,
\item descriptive records.
\end{itemize}
\end{example}

This list is illustrative; the analysis is not restricted to these forms.

These categories are not introduced as a fixed ontology.
They serve to distinguish the structural forms over which identity conditions
may be defined. Regime profiles refine these categories further
based on transformation behavior.


%--------------------------------------------
\subsection{Transformation Semantics}
\label{subsec:transformation_semantics}

This subsection specifies the interpretation
of the transformation families in $\mathbb{F}$.

\begin{definition}[Transformation semantics]
\label{def:transformation_semantics}
Each transformation family in $\mathbb{F}$ denotes a class of admissible
operations on representations with the following structural interpretation:
\begin{itemize}
\item $\mathrm{RE}$ (re-expression): change of representation format preserving content,
\item $\mathrm{AN}$ (annotation): addition of non-identity-bearing information,
\item $\mathrm{RF}$ (refinement): increase in descriptive or structural detail,
\item $\mathrm{AD}$ (aggregation/decomposition): restructuring into coarser or finer components,
\item $\mathrm{RC}$ (re-contextualization): change in applicability context,
\item $\mathrm{RA}$ (reassignment): change in association or bearer,
\item $\mathrm{SU}$ (substitution): replacement of the identity carrier,
\item $\mathrm{BF}$ (branch/fork): divergence into multiple continuations,
\item $\mathrm{PV}$ (provenance extension): addition of historical or derivational information,
\item $\mathrm{SE}$ (state evolution): change in state of an enduring referent over time.
\end{itemize}
\end{definition}


These interpretations are structural rather than interpretive.
They do not encode causal or normative commitments,
and are intended to apply uniformly across domains.
Each transformation family is defined extension-relatively
and is evaluated only through its effect on
equivalence relations induced by regime profiles.



Transformation families parameterize identity behavior.
A regime profile assigns to each applicable transformation
a classification determining whether identity is preserved,
broken, or ignored.
Distinct profiles therefore correspond
to distinct invariance conditions
over a fixed transformation basis.


%--------------------------------------------
\subsection{Profiles as Transformation Constraints}
\label{subsec:profiles_as_constraints}

This subsection characterizes how regime profiles
constrain admissible relations via induced identity.

\begin{definition}[Profile-governed relation]
\label{def:profile_governed_relation}
A relation $\rho \subseteq \mathcal{R} \times \mathcal{R}$ is profile-governed
if its admissibility is determined solely by the equivalence
classes $[x]_P$ of its relata under a fixed regime profile $P$.
\end{definition}


Preservation of a profile-governed relation under transformation
is equivalent to preservation of admissibility.
Relations depending on role, institutional context, or higher-order grouping
are not profile-governed and fall outside the scope of the present analysis.



A regime profile specifies how identity behaves under the transformation basis.
Different profiles correspond to distinct invariance conditions
over the same transformation families,
inducing distinct equivalence relations on $\mathcal{R}$.



Profile-governed relations depend only on equivalence classes.
Accordingly, all admissible structure is invariant under replacement
of representatives within a class, and representation proceeds
at the level of quotient structure.



The derivation of regime profiles proceeds by identifying coherent
transformation-classification patterns required for stable reference,
and separating those that cannot be realized by a single profile.
When no single classification assignment for a transformation family
is consistent across all admissible cases,
refinement is required.



Extension-stable and extension-invariant are used synonymously.

%--------------------------------------------
\subsection{Canonical Regime Classification Tables}
\label{subsec:canonical_classifications}

This subsection records the transformation classification for the three
regime profiles that do not exhibit split pressure under the transformation
basis $\mathbb{F}$.
Their classifications are stated here to make this paper self-contained
and to anchor the cross-family non-collapse arguments of
Section~\ref{subsec:pairwise_noncollapse}.

\subsubsection{OBL: Obligation-bearing entity}

\begin{itemize}
\item $\mathrm{RE}$: IGN
\item $\mathrm{AN}$: IGN
\item $\mathrm{RF}$: PRS
\item $\mathrm{AD}$: PRS
\item $\mathrm{RC}$: IGN
\item $\mathrm{RA}$: IGN
\item $\mathrm{SU}$: BRK
\item $\mathrm{BF}$: IGN
\item $\mathrm{PV}$: IGN
\item $\mathrm{SE}$: IGN
\end{itemize}

Identity is determined by persistence as a responsible party.
Obligation-bearing entities do not undergo state evolution in the sense
relevant to $\mathrm{SE}$: their identity is not individuated
by internal state transitions but by referential continuity as a bearer.
Branching is identity-ignoring: the obligation-bearing entity
persists independently of how its representations are forked.

\subsubsection{OCC: Time-indexed occurrence}

\begin{itemize}
\item $\mathrm{RE}$: IGN
\item $\mathrm{AN}$: IGN
\item $\mathrm{RF}$: PRS
\item $\mathrm{AD}$: BRK \quad
      (decomposition produces sub-occurrences with distinct temporal
      individuation; the original occurrence's identity is not preserved)
\item $\mathrm{RC}$: IGN
\item $\mathrm{RA}$: N/A \quad ($\mathrm{Appl}_{\mathrm{OCC}}(\mathrm{RA}) = 0$)
\item $\mathrm{SU}$: BRK
\item $\mathrm{BF}$: N/A \quad ($\mathrm{Appl}_{\mathrm{OCC}}(\mathrm{BF}) = 0$)
\item $\mathrm{PV}$: IGN
\item $\mathrm{SE}$: PRS
\end{itemize}

Identity is determined by temporal realization and provenance.
State evolution is identity-preserving: an occurrence may have
stages or phases that constitute the same event.
Decomposition is identity-breaking: it produces sub-occurrences
with distinct temporal individuation, and the original occurrence's
identity is not preserved.
Branching and reassignment are inapplicable because a time-indexed
occurrence is fixed by its realization event and admits no forking
or bearer reassignment while preserving that individuation.

\subsubsection{REC: Descriptive record}

\begin{itemize}
\item $\mathrm{RE}$: IGN
\item $\mathrm{AN}$: PRS
\item $\mathrm{RF}$: PRS
\item $\mathrm{AD}$: PRS
\item $\mathrm{RC}$: IGN
\item $\mathrm{RA}$: IGN
\item $\mathrm{SU}$: BRK
\item $\mathrm{BF}$: PRS
\item $\mathrm{PV}$: PRS
\item $\mathrm{SE}$: IGN
\end{itemize}

Identity is determined by persistence as a descriptive referent
without causal or normative commitment.
Annotation, provenance extension, and branching are identity-preserving:
a descriptive record persists across enrichment, derivation, and copying.
State evolution is identity-ignoring: descriptive records do not
have internal states in the sense relevant to $\mathrm{SE}$.

\medskip
\noindent\textbf{Canonical classification summary.}

\[
\begin{array}{@{}c|ccc@{}}
 & \mathrm{OBL} & \mathrm{OCC} & \mathrm{REC} \\
\hline
\mathrm{RE} & \mathrm{IGN} & \mathrm{IGN} & \mathrm{IGN} \\
\mathrm{AN} & \mathrm{IGN} & \mathrm{IGN} & \mathrm{PRS} \\
\mathrm{RF} & \mathrm{PRS} & \mathrm{PRS} & \mathrm{PRS} \\
\mathrm{AD} & \mathrm{PRS} & \mathrm{BRK} & \mathrm{PRS} \\
\mathrm{RC} & \mathrm{IGN} & \mathrm{IGN} & \mathrm{IGN} \\
\mathrm{RA} & \mathrm{IGN} & \mathrm{N/A} & \mathrm{IGN} \\
\mathrm{SU} & \mathrm{BRK} & \mathrm{BRK} & \mathrm{BRK} \\
\mathrm{BF} & \mathrm{IGN} & \mathrm{N/A} & \mathrm{PRS} \\
\mathrm{PV} & \mathrm{IGN} & \mathrm{IGN} & \mathrm{PRS} \\
\mathrm{SE} & \mathrm{IGN} & \mathrm{PRS} & \mathrm{IGN} \\
\end{array}
\]
